\section{Related Works}
\label{sec:related_works}

\paragraph{Generalization Problems in Robotic Manipulation.}
Generalization has been a central problem in robotic manipulation, and existing studies mainly evaluate it from three perspectives: scene-level, category-level, and cross-embodiment generalization.
Scene-level generalization focuses on whether a policy can remain robust when the surrounding visual or physical contexts vary, such as changes in object layouts, lighting conditions, or camera viewpoints~\cite{goyal2023rvt, shridhar2023perceiver, shridhar2022cliport, nair2023r3m}. 
Furthermore, category-level generalization studies whether robots can transfer manipulation skills to unseen objects that belong to the same semantic category while differing in geometry, size, or texture~\cite{gao2021kpam, mo2021where2act, geng2023gapartnet}.
More recently, cross-embodiment generalization has attracted increasing attention, which aims to train a unified policy that can transfer across different robot platforms, such as robotic arms and humanoids~\cite{o2024open, team2024octo, intelligence2025pi}.
In this work, we go beyond generalization over visual appearance and robot embodiment and focus on \textit{functional generalization}, in which an agent is required to accomplish the same function using diverse, unseen tools.


\paragraph{Improving Generalization in Robot Policy Learning.}
To improve manipulation generalization, the most straightforward way is scaling up action-labeled robotic data and training policies end-to-end~\cite{zitkovich2023rt, o2024open, team2024octo, khazatsky2024droid}. 
However, collecting large-scale robotic data is time-consuming and labor-intensive, especially for complicated tool-use manipulations. 
Reducing the dependence on action labels, some works learn transferable motion priors from large-scale action-free video data~\cite{nair2023r3m, wang2023mimicplay, wu2024unleashing,ma2026dit4dit, pai2025mimic, an2026feedback, hou2026world}. 
For example, MimicPlay~\cite{wang2023mimicplay} learns high-level latent plans from human play videos and grounds them into robot actions with a low-level visuomotor policy. 
Although videos contain rich motion and interaction cues, dense future frames do not explicitly reveal a compact abstraction of function-relevant intent. As a result, effectively learning functional priors from videos often requires extremely extensive pretraining.
Another line of work improves action grounding with more structured intermediate representations, such as interaction heatmaps~\cite{mo2021where2act, bahl2023affordances, wu2025afforddp, li2026compassad}, object-centric representations~\cite{zhulearning, chen2025tool, chapin2025object}, or keypoint trajectories~\cite{gao2021kpam, huang2025rekep, turpin2021gift, wen2023any}. 
For instance, ATM~\cite{wen2023any} predicts future trajectories of arbitrary points from video and uses them to guide a low-level policy for action generation. 
However, these methods mainly focus on standard manipulation tasks, such as object pick-up, while overlooking tool-use scenarios that require capturing more abstract functional intent, which are systematically investigated in our work. 
% through the lens of different intermediate representations.
% However, these methods mainly focus on standard manipulation tasks, such as object pick up, while overlooking tool-use scenarios that require understanding more abstract functional intent. 
% In this work, we systematically investigate the effect of different intermediate representations on functional generalization and identify that keypoint trajectories provide the best trade-off between visual predictability and action groundability.

%\paragraph{Intermediate Representations for Tool-use Manipulation.}
% Instead of directly predicting low-level actions from raw observations, many manipulation methods introduce intermediate representations to provide more structured guidance. Affordance-based methods identify task-relevant object regions for interaction, while keypoint-based methods represent objects or tools using sparse semantic points to support category-level manipulation and action grounding. Other works use human videos or predicted future observations as high-level plans, which can capture task intent before execution. These intermediate signals improve interpretability and transferability, but they often differ in visual predictability and action groundability. For contact-rich tool-use tasks, an effective representation should not only capture where the functional interaction occurs, but also describe how the tool should move over time. In this work, we systematically study several candidate representations, including affordance images, human video prompts, and 2D keypoint trajectories, and show that keypoint trajectories provide a better trade-off between functional planning and executable action grounding.

% Improving the generalization capability of robotic policies is a central challenge for deploying robots in open-world environments. Existing methods mainly address this problem from three perspectives.

% \paragraph{Scaling up Robotic Data.}
% One direct way to improve generalizable planning is to scale up robotic demonstrations and train policy models end-to-end to predict actions from observations~\cite{}. For example, Open X-Embodiment~\cite{} aggregates large-scale robot trajectories from many laboratories and enables RT-X to learn transferable skills across tasks and embodiments. Similarly, Octo~\cite{} trains a generalist policy on heterogeneous robot datasets and shows strong multi-task manipulation ability. However, collecting large-scale robot data is extremely labor-intensive and time-consuming, and the resulting policies may still lack explicit functional abstraction beyond the demonstrated objects and scenes.

% \paragraph{Pretraining on Human Videos.}
% Another direction is to pretrain policy models on action-free human videos to learn motion priors, and then transfer these priors to robot action generation through fine-tuning on a small amount of robotic demonstrations. For instance, R3M~\cite{} learns reusable visual representations from large-scale human videos and transfers them to downstream robot manipulation tasks. MimicPlay~\cite{} further extracts high-level plans from human play videos and grounds them with a low-level robot policy. However, without explicit action or physical supervision, it is difficult for models to learn the precise dynamics required for motor control, which may require extremely large-scale video pretraining.

% \paragraph{Exploiting Foundation Models.}
% Another stream of works~\cite{} uses pretrained foundation models to extract generalizable affordance or planning signals, which serve as intermediate bridges between seen and unseen demonstrations. For example, RT-H~\cite{} leverages vision-language reasoning to predict action hierarchies, allowing high-level semantic plans to guide low-level robot execution. Large Video Planner further takes a video foundation model as a high-level planner to predict future visual plans, which are then grounded into robot actions by a low-level policy. However, using foundation models during inference can introduce significant latency, and the final policy performance is strongly limited by the accuracy of the predicted intermediate signals.

% \paragraph{Summary.}
% Beyond data-centric strategies, recent methods~\cite{} exploit pretrained foundation models to extract generalizable affordance or planning signals, which serve as intermediate bridges between seen and unseen demonstrations. The goal is to enable a policy to transfer the same functional behavior, such as hitting, to novel tools with different shapes, materials, and interaction regions. To this end, we propose an efficient System-2 plus System-1 architecture, where System-2 predicts generalizable intermediate functional plans and System-1 grounds these plans into executable robot actions. This design aims to bridge high-level functional reasoning and low-level action execution, thereby improving policy generalization beyond object-specific imitation.
